

\section{Introducción}
\label{inicio_cap1}

El presente Manual de Operaciones (MO) establece la estructura organizacional, las políticas, los procedimientos operacionales, las responsabilidades y los controles adoptados por Skyfilms S.A.S. para la planificación, ejecución y supervisión de las operaciones con Sistemas de Aeronaves No Tripuladas (UAS) desarrolladas bajo la categoría específica, de conformidad con lo establecido en el Reglamento Aeronáutico de Colombia (RAC) 100 y demás disposiciones emitidas por la Unidad Administrativa Especial de Aeronáutica Civil (Aerocivil).

Las operaciones cubiertas por este Manual son de \textbf{simple captura de imágenes o datos} con fines audiovisuales, bajo la modalidad de contacto visual \textbf{VLOS}, con aeronaves DJI Mavic 4 Pro Creator Combo y BetaFPV Pavo20 Pro II, e incluyen los vuelos especiales de \textbf{vuelo nocturno} y \textbf{vuelo en zona urbana}. Quedan excluidos de forma expresa EVLOS, BVLOS, el vuelo autónomo o automático y los vuelos en espacios cerrados o confinados.

El cumplimiento de las disposiciones contenidas en este manual es obligatorio para todo el personal que participe en las operaciones UAS de Skyfilms, quienes deberán conocer, aplicar y mantener los procedimientos aquí establecidos durante el desarrollo de sus funciones.


\textbf{Nota:} Este Manual de Operaciones establece los procedimientos operacionales y administrativos de la organización. Los aspectos específicos relacionados con el Sistema de Gestión de Seguridad Operacional (SMS) y el Control de Mantenimiento se desarrollan en sus respectivos manuales, manteniendo coherencia, trazabilidad y control documental entre todos los documentos del sistema.


\section{Objetivo}

El presente Manual de Operaciones tiene por objeto establecer los lineamientos, políticas, procedimientos y responsabilidades que regulan la planificación, ejecución, supervisión y control de las operaciones con Sistemas de Aeronaves No Tripuladas (UAS) desarrolladas por Skyfilms S.A.S., con el fin de garantizar una operación segura, estandarizada y conforme a los requisitos establecidos en el Reglamento Aeronáutico de Colombia (RAC) 100 y demás disposiciones emitidas por la Aerocivil.



\section{Alcance}
El presente Manual de Operaciones aplica a todas las actividades de planificación, ejecución, supervisión y control de las operaciones realizadas por Skyfilms S.A.S. como explotador de Sistemas de Aeronaves No Tripuladas (UAS) certificado en la Categoría Específica, para la prestación de servicios de captura de imágenes y datos con fines audiovisuales, cinematográficos, publicitarios y comerciales.

Las operaciones se desarrollarán bajo la modalidad de \textbf{Visibilidad en Línea de Vista (VLOS)}, incluyendo autorizaciones operacionales para \textbf{Vuelo Nocturno}, dentro del territorio de la República de Colombia. Todas las actividades se ejecutarán bajo las condiciones, limitaciones y autorizaciones establecidas en el RAC 100, las OpSpecs aprobadas por la Aerocivil y los escenarios operacionales autorizados, utilizando únicamente las aeronaves, equipos tecnológicos y sistemas declarados por Skyfilms ante la autoridad aeronáutica.

Las disposiciones contenidas en este manual son de obligatorio cumplimiento para todo el personal que participe directa o indirectamente en las operaciones UAS de Skyfilms, incluyendo el Gerente General, Jefe de Pilotos UAS, Gerente SMS (en adelante, Gerente SMS), pilotos remotos, observadores visuales y demás personal de apoyo operacional.

\textbf{Tipo de operación:} simple captura de imágenes o datos. \textbf{Contacto visual:} VLOS. \textbf{Aeronaves:} DJI Mavic 4 Pro Creator Combo (distancia máxima de 750 m) y BetaFPV Pavo20 Pro II (distancia máxima prevista de 250 m, siempre con observador). \textbf{Vuelos especiales:} nocturno y zona urbana. \textbf{Altura máxima:} 400 ft AGL. Estas condiciones coinciden con las OpSpecs propuestas y con la Declaración de Intención de Solicitud (DIS).

El presente Manual \textbf{no contempla y excluye expresamente} operaciones en modalidad de visibilidad extendida (EVLOS), operaciones más allá de la línea de vista visual (BVLOS), transporte de carga (\textit{Drone Delivery}), aspersión, lanzamiento de objetos, operaciones autónomas sin supervisión, vuelos en espacios cerrados o confinados, vuelos de demostración comercial o de capacidad tecnológica, vuelos experimentales, o transporte de personas.


\section{Definiciones y Abreviaturas}

Las definiciones de los términos técnicos y operacionales empleados en el presente Manual de Operaciones se encuentran contenidas en las secciones Definiciones y Abreviaturas y Acrónimos, ubicadas después de la Tabla de Contenido del presente documento, las cuales hacen parte integral del mismo.




\section{Vigencia y Actualización}

El presente Manual de Operaciones entra en vigencia a partir de su aprobación por la Unidad Administrativa Especial de Aeronáutica Civil (Aerocivil) y permanecerá vigente mientras Skyfilms S.A.S. mantenga su certificado como explotador UAS, salvo que sea reemplazado por una versión posterior debidamente aprobada.

Toda modificación derivada de cambios normativos, organizacionales, técnicos u operacionales deberá incorporarse mediante el procedimiento de control documental establecido por Skyfilms. Cada revisión o enmienda será registrada en el Control de Revisiones del manual, comunicada al personal involucrado y, cuando corresponda, notificada o sometida a aprobación de la Aerocivil.

La emisión de una nueva versión del Manual de Operaciones dejará sin efecto las versiones anteriores, las cuales serán archivadas conforme al procedimiento de conservación documental de Skyfilms, garantizando su trazabilidad y consulta histórica cuando sea requerida.



\section{Información General del Explotador}


\begin{table}[H]
\centering
\renewcommand{\arraystretch}{1.35}
\begin{tabularx}{\textwidth}{|>{\bfseries}p{5cm}|X|}
\hline
Razón social & Skyfilms S.A.S. \\ \hline
Nombre comercial & Skyfilms \\ \hline
NIT & 902039520-1 \\ \hline
Representante legal & Andrés Felipe Pérez Sierra \\ \hline
Cargo & Gerente General \\ \hline
Base principal de operaciones & Ibagué - Tolima, Colombia \\ \hline
Cobertura operacional & República de Colombia \\ \hline
Actividad principal & Captura de imágenes y datos mediante UAS \\ \hline
Categoría operacional & Categoría Específica \\ \hline
Correo electrónico & skyfilmsfpv@gmail.com \\ \hline
Teléfono & +57 301 2646892 \\ \hline
Página web & www.skyfilms.com.co \\ \hline
Autoridad aeronáutica & Unidad Administrativa Especial de Aeronáutica Civil \\ \hline
\end{tabularx}
\caption{Información general del explotador UAS.}
\end{table}





Skyfilms S.A.S. es una organización especializada en la operación de Sistemas de Aeronaves No Tripuladas (UAS) para la captura de imágenes y datos, certificada como explotador UAS en la Categoría Específica y orientada a la prestación de servicios para la industria audiovisual, cinematográfica, comercial, publicitaria, inspección visual y demás actividades autorizadas por la Aerocivil.

Las operaciones se desarrollan bajo un enfoque de gestión del riesgo, seguridad operacional y cumplimiento normativo, de conformidad con el Reglamento Aeronáutico de Colombia (RAC 100), las Especificaciones Operacionales (OpSpecs) vigentes y los procedimientos establecidos en el presente Manual de Operaciones.







\section{Referencias Normativas y Documentales}

Los documentos relacionados a continuación constituyen las principales referencias normativas, técnicas y documentales empleadas para la elaboración, aplicación y actualización del presente Manual de Operaciones. Salvo indicación en contrario, Skyfilms utilizará siempre la versión vigente publicada por la autoridad competente o por el fabricante correspondiente.

\begin{table}[H]
\centering
\renewcommand{\arraystretch}{1.3}
\small
\begin{tabularx}{\textwidth}{|p{4cm}|p{3cm}|X|}
\hline
\textbf{Documento} & \textbf{Entidad} & \textbf{Aplicación} \\ \hline

Convenio sobre Aviación Civil Internacional (Anexos aplicables) &
OACI &
Marco internacional de la aviación civil. \\ \hline

Decreto 1294 de 2021 &
Gobierno de Colombia &
Reglamentación nacional aplicable a los sistemas de aeronaves no tripuladas. \\ \hline

RAC 100 &
Aerocivil &
Reglamento principal para la operación de Sistemas de Aeronaves No Tripuladas (UAS). \\ \hline

RAC 91 &
Aerocivil &
Reglas generales de operación aérea aplicables cuando corresponda. \\ \hline

AIP Colombia &
Aerocivil &
Información aeronáutica permanente para la planificación y ejecución de vuelos. \\ \hline

NOTAM &
Aerocivil &
Información aeronáutica temporal que puede afectar la operación. \\ \hline

Directiva Vinculante MAUT-5.0-22-016 &
Aerocivil &
Criterios de aceptación y gestión del enlace C2. \\ \hline

Circular Informativa MAUT-5.0-22-010 &
Aerocivil &
Registro de UAS, explotadores y equipos tecnológicos. \\ \hline

Manual de Operaciones (MO) &
Skyfilms S.A.S. &
Documento principal para la gestión operacional del explotador. \\ \hline

Manual del Sistema de Gestión de la Seguridad Operacional (SMS) &
Skyfilms S.A.S. &
Gestión de peligros, riesgos y seguridad operacional. \\ \hline

Manual de Control de Mantenimiento (MCM) &
Skyfilms S.A.S. &
Control y mantenimiento de aeronaves y equipos. \\ \hline

Especificaciones Operacionales (OpSpecs) &
Aerocivil &
Autorizaciones, limitaciones y escenarios operacionales aprobados. \\ \hline

Manuales técnicos del fabricante &
DJI, GEPRC, iFlight, Radiomaster, etc. &
Operación, mantenimiento y limitaciones de los equipos utilizados por Skyfilms. \\ \hline

\end{tabularx}

\caption{Referencias normativas y documentales aplicables al Manual de Operaciones.}
\end{table}

Las modificaciones o actualizaciones de cualquiera de los documentos anteriormente relacionados deberán ser evaluadas por el Gerente General y el Jefe de Pilotos UAS con el fin de determinar su impacto sobre el presente Manual de Operaciones. Cuando dichas modificaciones afecten los procedimientos operacionales, Skyfilms realizará la correspondiente revisión o enmienda del manual conforme al procedimiento de control documental establecido.







\section{Control del Manual y Revisiones}

El presente Manual de Operaciones de \textbf{Skyfilms} es un documento controlado. Toda copia impresa o digital debe estar identificada con su número de versión y fecha de emisión, de forma que pueda garantizarse que el personal autorizado utiliza siempre la versión vigente.  

\subsection*{Responsable del Control}
El Gerente General es el responsable de autorizar la emisión de nuevas versiones del manual. El Gerente SMS será el encargado de velar porque todo el personal operativo disponga de la versión vigente y de retirar de circulación cualquier versión obsoleta.  

\subsection*{Registro de Revisiones}
Cada revisión del manual debe ser registrada en el cuadro de registro de cambios, indicando como mínimo:  
\begin{itemize}
  \item Número de versión.  
  \item Fecha de la revisión.  
  \item Responsable que autoriza la modificación.
  \item Descripción de los cambios introducidos.  
    
\end{itemize}

El historial de revisiones hace parte integrante del manual y debe estar disponible para consulta de la Aerocivil en caso de requerirse.


\section{Control del Manual y Enmiendas}

Las enmiendas al Manual de Operaciones son los cambios específicos que se realizan entre versiones completas, con el fin de ajustar procedimientos, corregir errores o actualizar información.  

\subsection*{Procedimiento de Enmiendas}
\begin{enumerate}
  \item Cualquier miembro del personal podrá reportar la necesidad de una enmienda al Gerente General o al Gerente SMS.  
  \item El Gerente SMS evaluará la pertinencia de la enmienda y la someterá a aprobación del Gerente General.  
  \item Una vez aprobada, la enmienda se incorporará en el documento maestro del manual, indicando claramente la sección modificada y la fecha de implementación.  
  \item Se actualizará el registro de control de cambios, dejando constancia de la modificación.  
  \item El Gerente SMS será responsable de comunicar la enmienda a todo el personal y asegurarse de que se reemplacen las versiones obsoletas.  
\end{enumerate}

\section{Procedimiento para Informar Revisiones y/o Enmiendas a la Aerocivil}

Skyfilms S.A.S. se compromete formalmente a informar a la Unidad Administrativa Especial de Aeronáutica Civil (Aerocivil) cualquier revisión, enmienda, ajuste, modificación o actualización que afecte el contenido del presente Manual de Operaciones, especialmente aquellas que impacten directamente la operación UAS certificada, el alcance autorizado o las Especificaciones Relativas a las Operaciones (OpSpecs).

Toda modificación podrá clasificarse como revisión general o enmienda parcial, y será gestionada bajo el siguiente procedimiento:

\begin{enumerate}
    \item El Gerente SMS consolidará los cambios y actualizará el cuadro de control de revisiones o enmiendas, incluyendo el listado de páginas afectadas.
    \item La versión revisada o a la que se le aplico la enmienda del Manual, junto con el informe de cambios, será remitida oficialmente a la Aerocivil a través de los canales establecidos.
\end{enumerate}

Se debe tener presente para cada una lo siguiente:\\

\textbf{Revisiones:}
\begin{itemize}
    \item La revisión no entrará en vigor ni será aplicada operacionalmente hasta contar con la aprobación expresa o acuse formal de recibo por parte de la Aerocivil.
\end{itemize}
    
\textbf{Enmiendas:}
\begin{itemize}
    \item Las enmiendas menores (ej. correcciones tipográficas, ajustes de redacción que no afecten procedimientos)  podrán implementarse inmediatamente; no obstante, deberán notificarse a la Aerocivil dentro de los cinco (5) días hábiles siguientes.
    \item Las enmiendas que modifiquen procedimientos operacionales, responsabilidades, limitaciones o condiciones técnicas requerirán aprobación previa de la Aerocivil antes de su aplicación.
\end{itemize}

En ningún caso se ejecutarán operaciones UAS bajo versiones del Manual que no hayan sido debidamente notificadas o aprobadas cuando corresponda.\\

El Gerente General será responsable de autorizar y firmar la comunicación oficial dirigida a la Aerocivil, garantizando su envío oportuno y la conservación del soporte documental correspondiente, incluyendo número de radicación, fecha y evidencia de recepción.\\

Toda comunicación remitida será archivada digitalmente en el repositorio documental institucional, asegurando trazabilidad, control de versiones y preservación del historial documental.

Además de las revisiones del Manual, Skyfilms informa a la Aerocivil, dentro de los quince (15) días siguientes a su ocurrencia, todo cambio en su información de registro: aeronaves, equipos tecnológicos, pilotos, organización, Jefe de Pilotos UAS e instalaciones (RAC 100.505 d). El Gerente General también actualiza una vez al año la información entregada a la Aerocivil como explotador certificado (RAC 100.510 a-18), y el Gerente SMS verifica que los datos del registro coincidan con los que respaldan cada solicitud de autorización de vuelo.







\section{Control y Procedimiento para Distribución del Manual}

El Manual de Operaciones de Skyfilms es un documento controlado cuya distribución al interior de la compañía está sujeta a procedimientos formales para garantizar que el personal autorizado disponga únicamente de la versión vigente.La distribucion de este manual se realizara a todo el personal al momento de su vinculacion a la organizacion y/o cada vez que se produzcan modificaciones de carácter normativo, técnico u operacional.

\subsection*{Responsables}
El Gerente General es la autoridad responsable de autorizar la distribución de nuevas versiones del manual. El Gerente SMS se encargará de coordinar la entrega de los ejemplares, ya sea en formato digital o impreso, y de mantener actualizado el registro de distribución.

\subsection*{Modalidades de Distribución}
La distribución del manual se realizará en dos modalidades:
\begin{itemize}
    \item \textbf{Formato digital:} Una copia en PDF será enviada al Jefe de Pilotos, piloto/s, observador/es, Técnico UAS y demás personal autorizado. Esta copia será de solo lectura y estará identificada con el número de versión y fecha de emisión.
    \item \textbf{Formato impreso:} Se conservará al menos un ejemplar físico en las oficinas de Skyfilms, disponible para consulta del personal y de la Aerocivil en caso de auditorías o inspecciones.
\end{itemize}

\subsection*{Procedimiento y momento de la distribución}
\begin{enumerate}
    \item El Gerente General aprueba la nueva versión y la autoriza para su distribución.
    \item El Gerente SMS entrega la copia a cada receptor dentro de los cinco (5) días hábiles siguientes a la aprobación y retira las copias obsoletas.
    \item Cada receptor firma el \textbf{Acuse de recibo del Manual} (F-015) y devuelve el original al Gerente SMS.
    \item El Gerente SMS actualiza el cuadro de control de distribución y archiva el acuse en el Sistema de Gestión Documental.
\end{enumerate}
El Jefe de Pilotos verifica que todo el personal operativo haya recibido la versión vigente antes de la siguiente operación. El Gerente SMS es el único responsable de la distribución.

\subsection*{Control de Distribución}
Cada entrega de copia del manual quedará registrada en el cuadro de control de distribución, con los siguientes datos mínimos:
\begin{itemize}
    \item Nombre del receptor.  
    \item Cargo o rol dentro de la compañía.  
    \item Número de versión entregada.  
    \item Fecha de entrega y número del acuse de recibo (F-015).
\end{itemize}

\subsection*{Limitaciones}
Está prohibida la reproducción no autorizada del manual, parcial o total. Las versiones anteriores deben ser retiradas inmediatamente de circulación al momento de distribuir una nueva versión. El personal que conserve versiones obsoletas será responsable de su destrucción segura o entrega al Gerente SMS para su archivo.


\section*{Persona Responsable de la Actualización del Manual}
El \textbf{Gerente del Sistema de Gestión de la Seguridad Operacional (SMS)} será el responsable directo de mantener actualizado el Manual de Operaciones (MO) de Skyfilms y todas sus partes.  

Sus funciones específicas incluyen:  
\begin{itemize}
    \item Verificar periódicamente los cambios normativos en los Reglamentos Aeronáuticos de Colombia (RAC) y demás circulares emitidas por la Aerocivil, aplicables a las operaciones UAS.  
    \item Coordinar con el Gerente General y el Jefe de Pilotos la incorporación de revisiones o enmiendas al MO.  
    \item Garantizar que todas las copias digitales e impresas del manual correspondan a la última versión aprobada.  
    \item Mantener actualizado el \textit{Registro de Revisiones y Enmiendas}, así como el \textit{Listado de Páginas Afectadas (LPA)}.  
    \item Comunicar a todo el personal autorizado la entrada en vigor de las revisiones o enmiendas, asegurando su comprensión y aplicación en las operaciones.  
\end{itemize}




\section{Sistema de Medidas Adoptado}

Con el fin de garantizar la estandarización de la información técnica y operativa, Skyfilms adopta oficialmente el Sistema Internacional de Unidades (SI) para la totalidad de sus actividades y registros operacionales, salvo en aquellos parámetros aeronáuticos que por convención internacional se expresan en otras unidades.

\subsection*{Parámetros generales}
\begin{itemize}
    \item Longitud y distancia: metros (m) y kilómetros (km).  
    \item Masa y peso: kilogramos (kg).  
    \item Tiempo: Hora Universal Coordinada (UTC).  
    \item Temperatura: grados Celsius (°C).  
    \item Presión atmosférica: hectopascales (hPa), conforme QNH local..
    \item {Visibilidad:} kilómetros (km) o metros (m).
    
\end{itemize}

\subsection*{Parámetros aeronáuticos convencionales}
\begin{itemize}
    \item \textbf{Altitud y altura sobre el terreno:} pies (ft) sobre el nivel medio del mar (MSL) o sobre el terreno (AGL), según corresponda.  
    \item \textbf{Velocidad:} metros por segundo (m/s) y/o nudos (kt).  
    \item \textbf{Distancia de navegación:} millas náuticas (NM).  
\end{itemize}

\subsection*{Aplicación}
Todas las listas de verificación, reportes de vuelo, registros de mantenimiento y demás documentos oficiales de Skyfilms emplearán de forma consistente el sistema de medidas aquí descrito. Cualquier excepción deberá estar debidamente justificada en el manual correspondiente o en un anexo técnico.

\subsection*{Conversión de unidades}
Cuando un dato se exprese en una unidad aeronáutica y deba usarse en otra, se aplican estas equivalencias: 1 ft = 0,3048 m; 1 kt = 0,5144 m/s; 1 NM = 1\,852 m. Las alturas y altitudes se registran en pies, las distancias horizontales en metros o kilómetros, y la velocidad en m/s; el nudo solo se emplea al citar datos meteorológicos. Las especificaciones de fabricante en otras unidades se convierten al sistema de este Manual. El sistema de coordenadas adoptado es WGS 84.
